\chapter{$H\mathbb F_2$-synthetic homotopy and synthetic Adams objects}
\label{chap:synthetic-foundations}

This chapter introduces the synthetic objects and the language in which the imported Pstragowski and BHS statements can later be written.

\section{Topological operad language}

% SOURCE: Def/HigherAlgebra/Operad/{Arity,Topological,Algebra}; May,
% The Geometry of Iterated Loop Spaces, Definitions 1.1--1.2 and Lemma 1.4;
% What precisely are E-infinity ring spaces and E-infinity ring spectra?,
% Section 1, pp. 224--225. The ambient spaces here are all TopCat spaces.
\begin{definition}[Topological operads and Cartesian algebras]
  \label{def:topological-operad-foundation}
  \lean{KIP126.HigherAlgebra.Operad.TopologicalOperad,KIP126.HigherAlgebra.Operad.TopologicalOperad.Hom}
  \lean{KIP126.HigherAlgebra.Operad.TopologicalOperad.SigmaFree,KIP126.HigherAlgebra.Operad.TopologicalOperad.IsEInfinity,KIP126.HigherAlgebra.Operad.TopologicalOperad.IsReduced}
  \lean{KIP126.HigherAlgebra.Operad.TopologicalOperad.Algebra,KIP126.HigherAlgebra.Operad.TopologicalOperad.Algebra.Hom,KIP126.HigherAlgebra.Operad.TopologicalOperad.BundledAlgebra}
  \leanok
  A topological operad assigns a space $\mathcal O(I)$ to each finite
  labelled set $I$, a homeomorphism $e_*$ to each bijection of input sets,
  an identity $1\in\mathcal O(\{*\})$, and jointly continuous substitutions
  \[
    \gamma:\mathcal O(I)\times\prod_{i\in I}\mathcal O(J_i)
       \longrightarrow\mathcal O\Bigl(\coprod_{i\in I}J_i\Bigr).
  \]
  Relabelling preserves identities and composition. Left and right unit
  substitution are identities after the canonical singleton relabellings.
  Substitution is associative after the canonical bijection between the
  two dependent-sum label sets. Both equivariance laws are required: for
  $e:I\simeq I'$ and $e_i:J_i\simeq J'_i$,
  \[
    (\coprod e)_*\gamma(x,(y_{e(i)})_i)
       =\gamma(e_*x,(y_{i'})_{i'}),\qquad
    (\coprod_i e_i)_*\gamma(x,(y_i)_i)
       =\gamma(x,((e_i)_*y_i)_i).
  \]
  All laws include empty input sets. An operad morphism is continuous in
  every arity and preserves relabelling, identity, and substitution.

  In the chosen topological convention, $\mathcal O$ is $E_\infty$ when
  every $\mathcal O(I)$ is contractible and its permutation action is free.
  Contractibility means homotopy equivalence to a point. The separate
  reduced condition requires exactly one nullary operation; it is not
  imposed on every operad. A Cartesian $\mathcal O$-algebra on a space $X$
  has jointly continuous actions $\mathcal O(I)\times X^I\to X$ preserving
  the identity, substitution, and relabelling. Algebra morphisms preserve
  every such action, including the nullary ones.
\end{definition}

This definition supplies the topological operad language only. Its ambient
category is all topological spaces with ordinary finite products, without
May's additional compactly generated Hausdorff convention. Applying it to
synthetic spectra still needs enriched tensor actions, an applicable model
and its comparison with the chosen synthetic context, and derived mapping
spaces for the unit-preserving algebra structures. Those missing interfaces
include the meaning of the required uniqueness statement; the Cartesian
space-valued algebra definition does not complete checklist item a09.

\section{The category, spheres, and deformation element}

% SOURCE: KIP126/Main/Axiom/Literature/MainPaper/main.tex:755-758; Pstragowski, construction of Syn_E.
\begin{definition}[Synthetic context]
  \label{def:synthetic-context}
  \uses{def:stable-context,def:hf2-homology}
  \notready
  An $H\mathbb F_2$-synthetic context is a stable, presentable, symmetric
  monoidal category $\operatorname{Syn}_{H\mathbb F_2}$, together with a lax
  symmetric monoidal, filtered-colimit-preserving functor
  $\nu:\operatorname{Sp}\to\operatorname{Syn}_{H\mathbb F_2}$.  It has
  bigraded suspension, bigraded homotopy groups, cofibers, smash products, and
  the coherences required by the constructions below.
\end{definition}

% SOURCE: Pstragowski Definition 4.6; KIP126/Main/Axiom/Literature/MainPaper/main.tex:770-772.
\begin{definition}[Synthetic spheres]
  \label{def:synthetic-sphere}
  \uses{def:synthetic-context}
  \notready
  For integers $s,w$, define the synthetic sphere
  \[
    S^{s,w}=\Sigma^{s-w}\nu(S^w).
  \]
  Suspension by $S^{s,w}$ is written $\Sigma^{s,w}$ and the corresponding
  bigraded homotopy group is $\pi_{s,w}$.
\end{definition}

% SOURCE: Pstragowski Definition 4.27; KIP126/Main/Axiom/Literature/MainPaper/main.tex:774-782.
\begin{definition}[The deformation element and its quotients]
  \label{def:synthetic-lambda-quotient}
  \uses{def:synthetic-sphere}
  \notready
  The comparison map $S^{0,-1}\to S^{0,0}$ is denoted
  $\lambda\in\pi_{0,-1}S^{0,0}$.  For any synthetic spectrum $X$, multiplication
  by $\lambda^n$ is a natural map
  $\Sigma^{0,-n}X\to X$, and $X/\lambda^n$ is its cofiber.  In particular,
  $X/\lambda\simeq(S^{0,0}/\lambda)\mathbin{\wedge}X$.
\end{definition}

% SOURCE NOTE: KIP126/Main/Axiom/Literature/MainPaper/main.tex:784--786; Burklund--Xu Construction 7.7
% (KIP126/Main/Axiom/Literature/Sources/BurklundXu/paper.txt:3017--3039), using the filtered deformation
% machinery of Burklund--Hahn--Senger, Example C.15 and Appendix C.
\paragraph{Remark (higher $\lambda$-power quotient algebras).}
The claim in \texttt{KIP126/Main/Axiom/Literature/MainPaper/main.tex:784--786} is true: the locally filtered
synthetic category supplies a coherent tower of commutative algebra objects
$S^{0,0}\to\cdots\to S^{0,0}/\lambda^3\to S^{0,0}/\lambda^2\to
S^{0,0}/\lambda$.  Its precise source chain is Burklund--Xu Construction 7.7,
which invokes Burklund--Hahn--Senger Example C.15 and the surrounding filtered
deformation construction.  This is an expository provenance record, not a
formalization obligation: the project models the quotient maps and their
coherence below, while a complete $\mathrm{CAlg}$ tower and its higher
coherences remain outside the Project Boundary.

\section{Synthetic Adams objects}

% SOURCE: BHS Appendix A grading; KIP126/Main/Axiom/Literature/MainPaper/main.tex:799-811.
\begin{definition}[BHS-to-AIM regrading]
  \label{def:bhs-aim-regrading}
  \uses{def:synthetic-tridegree,def:reindexed-ss-morphism}
  \notready
  If BHS uses indices $(s,k,w_{\mathrm{BHS}})$, the paper uses
  \[
    (s,t,w_{\mathrm{AIM}})=(s,s+k,k+w_{\mathrm{BHS}}).
  \]
  Under this change a classical class in $(s,t)$ has synthetic degree
  $(s,t,t)$, $\lambda$ has degree $(0,0,-1)$, and differentials preserve the
  third coordinate.  The adapter includes proofs that page shapes and all
  displayed source and target degrees agree.
\end{definition}

% SOURCE: KIP126/Main/Axiom/Literature/MainPaper/main.tex:799-807, definition implicit in Theorem thm:rigid.
\begin{definition}[Synthetic Adams spectral sequence]
  \label{def:synthetic-adams-ss}
  \uses{def:synthetic-context,def:synthetic-tridegree,def:ss-elementwise-page-presentation,def:ss-einfty-presentation,def:multiplicative-ss,def:ss-convergence-detection-witness,def:unital-multiplication-data}
  \lean{KIP126.Synthetic.SpectralSequence.SyntheticAdamsSS}
  \lean{KIP126.Synthetic.SpectralSequence.SyntheticAdamsFamily}
  \lean{KIP126.Synthetic.SpectralSequence.SyntheticAdamsFamily.nu}
  \lean{KIP126.Synthetic.SpectralSequence.SyntheticAdamsFamily.nuQuotient}
  \lean{KIP126.Synthetic.SpectralSequence.SyntheticAdamsFamily.nuSphereIso}
  \lean{KIP126.Synthetic.SpectralSequence.SyntheticAdamsFamily.quotientProjection}
  \lean{KIP126.Synthetic.SpectralSequence.SyntheticAdamsFamily.deformationMap}
  \lean{KIP126.Synthetic.SpectralSequence.SyntheticAdamsConvergence}
  \notready
  For a synthetic spectrum in scope, its synthetic Adams spectral sequence is
  a trigraded spectral sequence
  \[
    {}^{\mathrm{syn}}E_r^{s,t,w}\Longrightarrow\pi_{t-s,w},
    \qquad d_r:(s,t,w)\mapsto(s+r,t+r-1,w),
  \]
  with finite-page elementwise cycles, boundaries, naturality, and an algebraic
  $E_\infty$ presentation.  A concrete convergence-and-detection witness is
  supplied separately whenever the paper identifies a synthetic abutment.  It has natural external smash pairings, and
  every such sequence is a module over the sequence for the synthetic sphere
  $S^{0,0}$.  The sphere sequence is internally multiplicative.  More
  generally, an internal multiplicative spectral sequence is obtained only
  for a synthetic spectrum equipped with chosen unital multiplication data;
  no internal product is asserted for an arbitrary synthetic spectrum.

  The internal family interface and its actual object bindings are defined:
  a single functor on the chosen synthetic category is evaluated at the
  chosen $\nu X$ and at the cofiber of the same $\lambda^n$ map.  Quotient
  projections and deformation maps are images of those actual morphisms.
  The explicit unit isomorphism of the same $\nu$ datum identifies its
  sphere image with the synthetic unit; the family transports this
  isomorphism to the internal sequences.
  The convergence datum uses the actual bigraded homotopy groups of that
  object.  Constructing the family, its pairings, and the needed convergence
  witnesses remains unfinished; these declarations do not choose a model or
  provide the literature results.
\end{definition}

\section{Normalized synthetic map data}

% SOURCE: KIP126/Main/Axiom/Literature/MainPaper/main.tex:929-961, Proposition prop:41561db2, its following
% torsion remark, and Notation not:fhat.
\begin{definition}[Chosen normalized synthetic map]
  \label{def:synthetic-hat-map}
  \uses{def:synthetic-context,def:synthetic-lambda-quotient,def:adams-filtration,def:cofiber-triangle}
  \notready
  For a selected occurrence of a map $f:X\to Y$ in a distinguished triangle,
  a triangle-compatible normalized synthetic map datum consists of the
  exponent $e(f)$, a chosen map $\widehat f$, its factorization equality, and
  its compatibility with one chosen synthetic lift of that triangle.  Put
  $e(f)=0$ when
  $\operatorname{AF}(f)=0$ and $e(f)=1$ when
  $\operatorname{AF}(f)>0$, including the infinite-filtration zero map.  A
  triangle-compatible normalized lift is the corresponding component
  \[
    \widehat f:\Sigma^{0,e(f)}\nu X\longrightarrow\nu Y,
    \qquad \nu f=\lambda^{e(f)}\widehat f,
  \]
  of that chosen lifted triangle after rotating back.  When $e(f)=0$
  this component is $\nu f$.  When $e(f)=1$, it is the particular component of
  the lifted triangle, not the map $\lambda^{k-1}\widetilde f$ formed from an
  arbitrary full lift.  This definition packages the required data but makes
  no existence claim; the located synthetic lift and triangle-lift results in
  the external-results chapter supply instances later.
\end{definition}

% SOURCE: KIP126/Main/Axiom/Literature/MainPaper/main.tex:929-971, the three rotations used in Proposition
% prop:41561db2 and Notation not:fhat; project disjoint packaging of the
% filtration hypotheses, including degenerate homology objects.
\begin{definition}[Normalized cofiber exactness case]
  \label{def:normalized-synthetic-exactness-case}
  \uses{def:synthetic-hat-map,def:hf2-homology,def:cofiber-triangle}
  \notready
  For a cofiber triangle
  $X\xrightarrow fY\xrightarrow gC(f)\xrightarrow h\Sigma X$, a normalized
  cofiber exactness case is one of the following tagged, disjoint data:
  \begin{enumerate}
    \item a zero case with $H_*f=0$ and
      $(e(f),e(g),e(h))=(1,0,0)$;
    \item a monomorphism case with $H_*f$ injective and
      $(e(f),e(g),e(h))=(0,0,1)$;
    \item an epimorphism case with $H_*f$ surjective and
      $(e(f),e(g),e(h))=(0,1,0)$.
  \end{enumerate}
  Each constructor also carries both the short exact homology sequence for the
  relevant rotation and one chosen lifted distinguished triangle datum for
  that rotation.  The latter field contains the triangle and its
  three compatible components; it is not merely an existence assertion.  The explicit
  $e$-pattern is part of the constructor, so degenerate maps that are both zero
  and mono or epi cannot be used with two contradictory shifts.  Equivalently,
  the data always includes $e(f)+e(g)+e(h)=1$ and records which summand is one.
\end{definition}
